% Gauss, Werke X/1 (1917), native leaves 29–32, printed 22–25.
\SourceLeaf{029}\NoSourceHead
\pdfbookmark[0]{Zur Theorie der Potenzreste}{potenzreste}
\HistoricTitle{[ZUR THEORIE DER POTENZRESTE.]}
\bigskip
\HistoricTitle{[QUADRATISCHE RESTE.]}
\SourceGroup{I}
\HistoricTitle{[ÜBER EINIGE SUMMEN.]}
\NumberedItem{1}
\HistoricCitation{[Aus Scheda Ae, Julius 1800, S. 28.]}
\begin{center}\textit{Theorema demonstrandum.}\end{center}
\[
\begin{array}{rl}
n&\text{numerus impar}\\
r&\text{radix aequationis }x^n-1=0
\end{array}
\]
\[
\left.\begin{array}{c}R\\N\end{array}\right\}
\text{indefinite }
\left\{\begin{array}{l}\text{residuum}\\\text{non residuum}\end{array}\right\}
\begin{array}{l}\text{ipsius }n,\ <n,\\\text{ad }n\text{ primum.}\end{array}
\]
\noindent Productum ex
\[
(r-r^{-1})(r^3-r^{-3})(r^5-r^{-5})\ldots(r^{n-2}-r^{n-2})
\]
\noindent evolvi in
\[
\Sigma r^R-\Sigma r^N.
\]
\noindent Demonstratum iam est illud productum evolvi vel in $\Sigma r^R-\Sigma r^N$ vel in $\Sigma r^N-\Sigma r^R$.
\ClosingRule
\clearpage

\SourceLeaf{030}\ArithmeticHead{23}{R}{ÜBER EINIGE SUMMEN.}
\NumberedItem{2}
\HistoricCitation{[Aus Scheda Ae, Julius 1800, S. 31.]}
\begin{center}\textit{Theorema novissimum pulcherrimum.}\end{center}
\[
\begin{array}{@{}l@{\quad}l@{}}
\text{Summae}&\begin{gathered}
\sin\frac0nP+\sin\frac1nP+\sin\frac4nP+\sin\frac9nP+\ldots+\sin\frac{(n-1)^2}{n}P\\[1ex]
\cos\frac0nP+\cos\frac1nP+\cos\frac4nP+\cos\frac9nP+\ldots+\cos\frac{(n-1)^2}{n}P
\end{gathered}
\end{array}
\]
\begingroup
\newcommand{\SumCell}[1]{\makebox[1.5em]{$#1$}}
\newsavebox{\SumBraceBox}
\sbox{\SumBraceBox}{$\left\{\vphantom{\begin{array}{c}\surd n\\\surd n\end{array}}\right.$}
\[
\begin{array}{r@{}l}
\text{fit}={}&\usebox{\SumBraceBox}\begin{array}{c|c|c|c|}
\SumCell{\surd n}&\SumCell{0}&\SumCell{0}&\SumCell{\surd n}\\
\SumCell{\surd n}&\SumCell{\surd n}&\SumCell{0}&\SumCell{0}
\end{array}\\[-.5ex]
\text{pro }n\equiv{}&\hspace{\wd\SumBraceBox}\underbrace{\begin{array}{|c|c|c|c|}\hline
\SumCell{0}&\SumCell{1}&\SumCell{2}&\SumCell{3}
\end{array}}_{\text{mod. }4}
\end{array}
\]
\endgroup
\ClosingRule
\NumberedItem{3}
\HistoricCitation{[Aus Scheda Ae, Julius 1800, S. 32.]}
\noindent Erstens $n\equiv1\ (\mathrm{mod.}\ 8)$
\[
\Sigma\cos\frac{RP}{n}-\Sigma\cos\frac{NP}{n}=+\surd n.
\]
\noindent Zweitens $n\equiv5\ (\mathrm{mod.}\ 8)$
\[
\Sigma\cos\frac{RP}{n}-\Sigma\cos\frac{NP}{n}=+\surd n.
\]
\noindent Drittens $n\equiv3\ (\mathrm{mod.}\ 8)$
\[
\Sigma\sin\frac{RP}{n}-\Sigma\sin\frac{NP}{n}=+\surd n.
\]
\noindent Viertens $n\equiv7\ (\mathrm{mod.}\ 8)$
\[
\Sigma\sin\frac{RP}{n}-\Sigma\sin\frac{NP}{n}=+\surd n.
\]
\ClosingRule
\clearpage

\SourceLeaf{031}\ArithmeticHead{24}{L}{ARITHMETIK.\quad NACHLASS.}
\NumberedItem{4}
\begin{center}\textsc{Gauss an Olbers.}\end{center}
\HistoricCitation{[\textsc{Wilhelm Olbers}, Sein Leben und seine Werke II, 1 (1900), S. 267 ff.]}
\begin{flushright}Braunschweig, 1805 September 3.\end{flushright}
\par\indent Ich hoffe, dass nur Ihre überhäuften Arbeiten, nicht aber Krankheit, oder sonst etwas Unangenehmes, Schuld sind, dass ich so lange mit keinem Briefe von Ihnen erfreut worden bin. Meine Beschäftigungen waren auch seit einiger Zeit nicht von der Art, dass sie für den Geometer, noch weniger die Ergebnisse, dass sie für den theilnehmenden Freund sonderlich Stoff zu Mittheilungen dargeboten hätten. Ich bin durch verschiedene Umstände — theils durch einige Briefe von \textsc{Le Blanc} in Paris, der meine \textit{Disqu. Arith.} mit wahrer Leidenschaft studirt, sich ganz mit ihnen vertraut gemacht und mir manche recht artige Kommunikationen darüber gemacht hat, theils durch die Anwesenheit eines Freundes, der jenes Werk jetzt gleichfalls studirt und sich öfters bei mir Raths erholt — theils auch durch eine Art von Überdruss oder wenigstens Ermüdung an dem todten mechanischen Kalkül verleitet worden, in diesem einmal eine Pause zu machen und meine geliebten arithmetischen Untersuchungen wieder vorzunehmen. Sie erinnern sich vielleicht noch von unsern Gesprächen in Bremen her, namentlich an dem schönen Nachmittage, den wir auf der Vahr zubrachten, dass ich schon seit längerer Zeit eine sehr beträchtliche Sammlung von Untersuchungen nicht sowohl im Pult, als in petto habe, die hinreichenden Stoff zu einem zweiten Bande der \textit{Disqu. Arith.} geben, und die, wenigstens meinem Urtheile nach, ebenso merkwürdig sind, als die im ersten enthaltenen. Sie erinnern sich aber auch vielleicht zu gleicher Zeit meiner Klagen über einen Satz, der theils schon an sich sehr interessant ist, theils einem sehr beträchtlichen Theile jener Untersuchungen als Grundlage oder als Schlussstein dient, den ich damals schon über 2 Jahr kannte, und der alle meine Bemühungen, einen genügenden Beweis zu finden, vereitelt hatte. Dieser Satz ist schon in meinen \textit{Disq.} pg. 636\begingroup\renewcommand{\thefootnote}{[*)}\footnote{Werke I, S. 442, 443.]}\textsuperscript{]}\endgroup\ oder vielmehr
\clearpage
\SourceLeaf{032}\ArithmeticHead{25}{R}{GAUSS AN OLBERS.}
\noindent nur ein specieller Fall davon, nämlich der, wo $n$ eine Primzahl ist, auf den sich übrigens hier die übrigen würden zurückführen lassen. Was da von »Quaecunque igitur radix etc.« bis »valde sunt memorabilia« steht, ist streng dort bewiesen, aber was folgt, nämlich die Bestimmung des Wurzelzeichens, ist es gerade, was mich immer gequält hat. Dieser Mangel hat mir alles Übrige, was ich fand, verleidet; und seit 4 Jahren wird selten eine Woche hingegangen sein, wo ich nicht einen oder den andern vergeblichen Versuch, diesen Knoten zu lösen, gemacht hätte — besonders lebhaft nun auch wieder in der letzten Zeit. Aber alles Brüten, alles Suchen ist umsonst gewesen, traurig habe ich jedesmal die Feder wieder niederlegen müssen. Endlich vor ein paar Tagen ist's gelungen — aber nicht meinem mühsamen Suchen, sondern bloss durch die Gnade Gottes möchte ich sagen. Wie der Blitz einschlägt, hat sich das Räthsel gelöst; ich selbst wäre nicht im Stande, den leitenden Faden zwischen dem, was ich vorher wusste, dem, womit ich die letzten Versuche gemacht hatte, — und dem, wodurch es gelang, nachzuweisen. Sonderbar genug erscheint die Lösung des Räthsels jetzt leichter als manches andere, was mich wohl nicht so viele Tage aufgehalten hat als dieses Jahre, und gewiss wird niemand, wenn ich diese Materie einst vortrage, von der langen Klemme, worin es mich gesetzt hat, eine Ahnung bekommen.
\par\medskip\noindent\makebox[\linewidth][s]{.\hfill.\hfill.\hfill.\hfill.\hfill.\hfill.\hfill.\hfill.\hfill.\hfill.\hfill.\hfill.\hfill.\hfill.\hfill.\hfill.\hfill.\hfill.\hfill.\hfill.\hfill.\hfill.\hfill.\hfill.\hfill.\hfill.}
\bigskip
\begin{center}\textbf{BEMERKUNG.}\end{center}
\begingroup\small
\par\indent Die vorstehenden Aufzeichnungen [1]—[3] gehören insofern in die Lehre von den quadratischen Resten als sie die ersten Vorstudien zu den Untersuchungen darstellen, die \textsc{Gauss} 1808 in der Abhandlung \textit{Summatio quarundam serierum singularium} (Werke II, S. 9 ff.) veröffentlicht hat. Für das Geschichtliche vergleiche man neben dem Briefe [4] noch die Tagebuchaufzeichnungen Nr. 118 vom Mai 1801 und Nr. 123 vom 30. August 1805 sowie den Artikel 18 des \textsc{Bachmann}schen Aufsatzes »Über \textsc{Gauss}' zahlentheoretische Arbeiten«.
\begin{flushright}\textsc{Schlesinger.}\end{flushright}
\ClosingRule\clearpage
\endgroup
